% from ega2-1-fr.tex
\oldpage[II]{177}
Si $Y$ est noethérien (non nécessairement affine) et $\mathcal{S}$ de type
fini, il y a un recouvrement fini de $Y$ par des ouverts affines noethériens
$U_i$, et on déduit alors de ce qui précède que pour tout $i$, il y a un
entier $n_i$ tel que
$\mathcal{S}_{n+1}|U_i=(\mathcal{S}_1|U_i)(\mathcal{S}_n|U_i)$ dès que
$n\geqslant n_i$~; le plus grand des $n_i$ vérifie donc
(\hyperref[II.8.7.5.1-fr]{8.7.5.1}).
\end{proof}

\begin{env}[8.7.6]
\label{II.8.7.6-fr}
Considérons maintenant le $C$-préschéma $Z$ obtenu en faisant \emph{éclater}
dans le cône affine $C$ le \emph{sous-préschéma des sommets}
$\varepsilon(Y)$~; par définition (\hyperref[II.8.1.3-fr]{8.1.3}) c'est donc
le préschéma
$\operatorname{Proj}(\bigoplus_{n\geqslant 0}\mathcal{S}_+^n)$~; l'injection
canonique
\begin{equation}
\label{II.8.7.6.1-fr}
  \iota:\bigoplus_{n\geqslant 0}\mathcal{S}_+^n
  \longrightarrow\mathcal{S}^{\natural}
\tag{8.7.6.1}
\end{equation}
définit (grâce à l'identification (\hyperref[II.8.7.3-fr]{8.7.3})) un
$C$-morphisme dominant canonique
\begin{equation}
\label{II.8.7.6.2-fr}
  G(\iota)\longrightarrow Z
\tag{8.7.6.2}
\end{equation}
où $G(\iota)$ est un ouvert de $C_X$ (\hyperref[II.3.5.1-fr]{3.5.1})~; on
notera qu'on peut avoir $G(\iota)\ne C_X$, comme le montre l'exemple où
$Y=\operatorname{Spec}(K)$, $K$ étant un corps, $\mathcal{S}=\widetilde{S}$
avec $S=K[\mathbf{y}]$, où $\mathbf{y}$ est une indéterminée \emph{de degré
2}~; si $R_n$ désigne l'ensemble $(S_+)^n$, considéré comme partie de
$S_{[n]}=S_n^{\natural}$, $S_+^{\natural}$ n'est pas la racine dans
$S^{\natural}$ de l'idéal engendré par la réunion des $R_n$ (cf.
(\hyperref[II.2.3.14-fr]{2.3.14})).
\end{env}

\begin{corollary}[8.7.7]
\label{II.8.7.7-fr}
Supposons qu'il existe $n_0>0$ tel que
\begin{equation}
\label{II.8.7.7.1-fr}
  \mathcal{S}_n=\mathcal{S}_1^n
  \qquad\text{pour }n\geqslant n_0.
\tag{8.7.7.1}
\end{equation}
Alors le morphisme canonique (\hyperref[II.8.7.6.2-fr]{8.7.6.2}) est partout
défini et est un isomorphisme $C_X\xrightarrow{\sim}Z$. Inversement, si cette
propriété a lieu et si on suppose de plus $Y$ noethérien et $\mathcal{S}$ de
type fini, il existe $n_0$ tel que l'on ait
(\hyperref[II.8.7.7.1-fr]{8.7.7.1}).
\end{corollary}

\begin{proof}
En effet, la première assertion est locale sur $Y$, et découle donc de
(\hyperref[II.8.2.14-fr]{8.2.14})~; il en est de même de la réciproque, en
raisonnant comme dans (\hyperref[II.8.7.5-fr]{8.7.5}).
\end{proof}

\begin{remark}[8.7.8]
\label{II.8.7.8-fr}
Comme la condition (\hyperref[II.8.7.7.1-fr]{8.7.7.1}) implique
(\hyperref[II.8.7.5.1-fr]{8.7.5.1}), on voit que lorsqu'elle est vérifiée,
non seulement $C_X$ s'identifie au préschéma obtenu en faisant éclater le
sommet (identifié à $Y$) du cône affine $C$, mais encore le sommet (identifié
à $X$) de $C_X$ s'identifie au sous-préschéma fermé image réciproque du
sommet $Y$ de $C$. En outre, l'hypothèse
(\hyperref[II.8.7.7.1-fr]{8.7.7.1}) implique que sur
$X=\operatorname{Proj}(\mathcal{S})$, les $\mathcal{O}_X$-Modules
$\mathcal{O}_X(n)$ sont inversibles (\hyperref[II.3.2.5-fr]{3.2.5} et
\hyperref[II.3.2.9-fr]{3.2.9}) et que l'on a
$\mathcal{O}_X(n)=\mathcal{L}^{\otimes n}$ avec
$\mathcal{L}=\mathcal{O}_X(1)$ (\hyperref[II.3.2.7-fr]{3.2.7} et
\hyperref[II.3.2.9-fr]{3.2.9})~; par définition
(\hyperref[II.8.6.1.1-fr]{8.6.1.1}), $C_X$ est donc le \emph{fibré vectoriel}
$\mathbf{V}(\mathcal{L})$ sur $X$, et son sommet la \emph{section nulle} de ce
fibré vectoriel.
\end{remark}

\subsection{Faisceaux amples et contractions}
\label{subsection:II.8.8-fr}

\begin{env}[8.8.1]
\label{II.8.8.1-fr}
Soient $Y$ un préschéma, $f:X\to Y$ un morphisme \emph{séparé et
quasi-compact}, $\mathcal{L}$ un $\mathcal{O}_X$-Module inversible \emph{ample
relativement à $f$}. Considérons la $\mathcal{O}_Y$-Algèbre graduée à degrés
positifs
\begin{equation}
\label{II.8.8.1.1-fr}
  \mathcal{S}=\mathcal{O}_Y\oplus
  \bigoplus_{n\geqslant 1}f_*(\mathcal{L}^{\otimes n})
\tag{8.8.1.1}
\end{equation}

